\documentclass[twoside,twocolumn,10pt]{article}

% -----------------------------
% Packages
% -----------------------------
\usepackage{fancyhdr}
\usepackage{geometry}
\usepackage{abstract}
\usepackage{graphicx}
\usepackage{titlesec}
\usepackage{ragged2e}
\usepackage{amssymb}
\usepackage{parskip}
\usepackage{amsmath}
\usepackage{newtxtext}
\usepackage{booktabs}
\usepackage{array}
\usepackage{balance}
\usepackage[backend=bibtex,style=ieee]{biblatex}
\addbibresource{references.bib}
\usepackage[font=footnotesize,labelfont=bf,justification=raggedright,format=plain]{caption}
\usepackage{float}

% -----------------------------
% Bibliography Heading
% -----------------------------
\defbibheading{bibliography}[\refname]{%
  \section{\MakeUppercase{REFERENCES}}}

% -----------------------------
% Custom Column Types
% -----------------------------
\newcolumntype{L}[1]{>{\raggedright\arraybackslash}p{#1}}
\newcolumntype{C}[1]{>{\centering\arraybackslash}p{#1}}
\newcolumntype{R}[1]{>{\raggedleft\arraybackslash}p{#1}}

% -----------------------------
% Caption Formatting
% -----------------------------
\captionsetup[table]{labelsep=period, justification=raggedright}
\captionsetup[figure]{labelsep=period, justification=raggedright}

% -----------------------------
% Page Geometry
% -----------------------------
\geometry{
    a4paper,
    left=0.8in,
    right=0.8in,
    top=1in,
    bottom=1in
}

\setlength{\columnsep}{15pt}
\setlength{\parindent}{0pt}
\setlength{\parskip}{0pt}
\hyphenpenalty=10000
\exhyphenpenalty=10000
\frenchspacing

% -----------------------------
% Abstract Formatting
% -----------------------------
\renewcommand{\abstractname}{\hspace{1.3em}\textbf{Abstract}}
\renewcommand{\abstractnamefont}{\normalfont\fontsize{10}{14}\bfseries\MakeUppercase}
\renewcommand{\absnamepos}{flushleft}
\renewcommand{\abstracttextfont}{\normalfont\fontsize{11}{13}\selectfont}

% -----------------------------
% Header / Footer Style
% -----------------------------
\setcounter{page}{1}
\pagestyle{fancy}
\fancyhf{}
\fancyhead[C]{\small RESOURCE-CONSTRAINED MULTIMODAL DISASTER ASSESSMENT}
\fancyhead[CO]{\small First Author et al. (2026)}
\fancyfoot[C]{\thepage}

% -----------------------------
% First Page Style
% -----------------------------
\fancypagestyle{firstpage}{
  \fancyhf{}
  \fancyhead[LO]{\small Volume X, Issue X, Year}
 \fancyhead[RO]{\small LGU J. Comput. Sci. \& IT}
  \fancyfoot[L]{\footnotesize Received: DD Month YYYY; Revised: DD Month YYYY; Accepted: DD Month YYYY\\\textsuperscript{*}These authors contributed equally to this work. Corresponding Author: First Author \textless author1@example.com\textgreater}
  \fancyfoot[R]{\footnotesize \textcopyright\ Year by the Author(s). Published by LGURJCSIT,\\ \hspace{0.5cm} under the CC-BY license}
  \renewcommand{\headrulewidth}{0.4pt}
  \renewcommand{\footrulewidth}{0.4pt}
}

% -----------------------------
% Title / Author
% -----------------------------
\title{\fontsize{16}{20}\selectfont\textbf{RESOURCE-CONSTRAINED MULTIMODAL DISASTER ASSESSMENT: A NOISE-RESILIENT FUSION ARCHITECTURE FOR LOW-RESOURCE AND HIGH-RESOURCE ENVIRONMENTS}}

\author{
    \fontsize{12}{14}\selectfont
    First Author\textsuperscript{1*}, Second Author\textsuperscript{2*}, Third Author\textsuperscript{3*}\\[1ex]
    \fontsize{12}{14}\selectfont
    \textsuperscript{1}Institution 1, City, Country \\
    \textsuperscript{2}Institution 2, City, Country \\
    \textsuperscript{3}Institution 3, City, Country
}

\date{}

% -----------------------------
% Section Formatting
% -----------------------------
\titleformat{\section}{\normalsize\bfseries\uppercase}{\thesection.}{1em}{}
\titleformat{\subsection}{\normalsize\bfseries\flushleft}{\thesubsection.}{1em}{}
\titleformat{\subsubsection}{\normalsize\bfseries\itshape\flushleft}{\thesubsubsection.}{1em}{}

% -----------------------------
% Document Start
% -----------------------------
\begin{document}

\twocolumn[
\begin{@twocolumnfalse}

\vspace{-0.7cm}
\noindent
\raggedright
{\fontsize{11}{13}\selectfont Original Article}
\hfill
{\fontsize{11}{13}\selectfont\footnotesize https://doi.org/xx.xxxx/xxxx}

\vspace{-0.5cm}

\maketitle
\thispagestyle{firstpage}

\section*{\hspace{-1.5mm}Abstract}
\fontsize{10}{12}\selectfont
\justifying
\noindent
Multimodal disaster assessment using Unmanned Aerial Vehicle (UAV) imagery and social media text is crucial for rapid response and resource allocation. However, deploying such models in resource-constrained or linguistically low-resource environments (e.g., Bengali text streams) poses significant challenges due to noisy data, extreme class imbalance, and hardware limitations. In this paper, we propose a noise-resilient dual-encoder late-fusion architecture specifically designed to tackle these challenges. Our model utilizes a lightweight vision backbone (Swin-Tiny) augmented with aggressive photometric distortions, coupled with a robust text encoder (BanglaBERT for local contexts, XLM-RoBERTa for global generalization). We introduce targeted Fast Gradient Method (FGM) adversarial training restricted strictly to word embeddings, preventing modality collapse. Furthermore, we deploy a Dynamic Pseudo-Labeling Strategy (DIPS) and an Optuna-optimized LightGBM stacker to calibrate log-bias shifts. Extensive benchmarking on the BanglaCalamityMMD and CrisisMMD v2.0 datasets demonstrates that our approach achieves a Macro F1 score of 0.745 in low-resource settings, outperforming state-of-the-art unimodal and baseline multimodal approaches.

\vspace{0.3cm}
\noindent\textbf{Keywords:}
Multimodal Disaster Assessment, Dual-Encoder Architecture, Fast Gradient Method, Pseudo-Labeling, Low-Resource Languages

\vspace{0.8cm}
\end{@twocolumnfalse}
]

% -----------------------------
% Main Text
% -----------------------------
\section{Introduction}
\fontsize{11}{13}\selectfont

Effective disaster management heavily relies on real-time situational awareness. With the proliferation of social media and drone surveillance, massive streams of multimodal data (text and imagery) are generated during catastrophic events. However, accurately classifying disaster severity from these streams is hindered by data noise, missing modalities, and extreme linguistic diversity. Specifically, low-resource languages like Bengali often suffer from insufficient annotated datasets and severe class imbalance, making standard transformer-based architectures prone to overfitting and routing collapse.

In this work, we present a robust, resource-constrained multimodal fusion architecture. We validate our framework through rigorous ablation on two datasets: the regional BanglaCalamityMMD dataset and the globally recognized CrisisMMD v2.0 dataset. 

Our main contributions are:
\begin{itemize}
    \item A late-fusion multimodal architecture utilizing Swin-Tiny and domain-specific BERT models optimized for low-resource environments.
    \item A targeted FGM adversarial training protocol that perturbs only the textual embedding space (\texttt{word\_embeddings}), safeguarding the highly augmented vision representations from adversarial degradation.
    \item The Dynamic Pseudo-Labeling Strategy (DIPS) coupled with a Level-2 LightGBM stacker for robust class calibration.
\end{itemize}

\section{Methodology}
\fontsize{11}{13}\selectfont

\subsection{Multimodal Dual-Encoder Late Fusion}
Our architecture independently processes image and text modalities to preserve spatial and semantic distinctiveness before late-stage projection. We employ \texttt{swin\_tiny\_patch4\_window7\_224} as the vision backbone. To counteract the scarcity of high-variance disaster visual data, we apply aggressive HSV (Hue, Saturation, Value) distortions. The text backbone uses \texttt{csebuetnlp/banglabert} for Phase 1 (Bengali) and \texttt{xlm-roberta-large} for Phase 2 (Multilingual). The pooled features from both encoders are concatenated, passed through a dropout layer ($p=0.3$), and projected to the target severity classes.

\subsection{Targeted Adversarial Training (FGM)}
Standard Fast Gradient Method (FGM) perturbs all trainable parameters based on gradient directions, which often disrupts the already highly-augmented visual manifolds in our pipeline. We introduce a \textit{Targeted FGM} module that strictly applies adversarial noise $r_{adv} = \epsilon \frac{g}{||g||_2}$ to the text encoder's \texttt{word\_embeddings}. This acts as a powerful regularizer for the text manifold, making the model highly resilient to social media typos and linguistic noise, without corrupting visual spatial coherence.

\subsection{FP16 Scaling and Safety}
To train the dual-encoder on resource-constrained hardware, we employ FP16 Mixed Precision. However, multimodal gradients frequently spike, leading to \texttt{NaN} scale factors. We implement a safe gradient scaler check: if $\text{scale}_{t} < \text{scale}_{t-1}$, the optimizer and learning rate scheduler steps are skipped, preventing gradient explosion and maintaining convergence stability.

\subsection{Dynamic Pseudo-Labeling Strategy (DIPS)}
To leverage unannotated disaster streams, we implement DIPS. DIPS tracks prediction distributions across multiple epochs on the test sets. Pseudo-labels are filtered dynamically based on a stringent confidence threshold ($\mu_{conf} \ge 0.85$) and low epistemic uncertainty ($\sigma_{conf} \le 0.15$). These verified labels are injected into a final fine-tuning phase.

\subsection{Optuna-Calibrated LightGBM Stacking}
Finally, to correct log-bias shifts induced by severe dataset imbalance (e.g., disproportionate `Normal' vs. `Severe' classes), we extract Out-Of-Fold (OOF) multimodal features and train a Level-2 LightGBM classifier. The hyperparameters are optimized via Optuna to maximize the Macro F1-score.

\section{Experiments}
\fontsize{11}{13}\selectfont

\subsection{Phase 1: BanglaCalamityMMD}
We conducted an extensive ablation study on the local BanglaCalamityMMD dataset to validate each architectural component. As shown in Table \ref{table1}, the base fusion model struggles with the complex noise distributions. Introducing Targeted FGM provides a significant boost, proving the efficacy of text-manifold regularization. The integration of DIPS and the LightGBM stacker pushes the architecture to state-of-the-art performance (Macro F1 = 0.745).

% -----------------------------
% Table 1
% -----------------------------
\begin{table}[h]
\centering
\caption{Ablation Study on BanglaCalamityMMD (Phase 1)}
\label{table1}
\renewcommand{\arraystretch}{1.2}
\setlength{\tabcolsep}{6pt}
\resizebox{\columnwidth}{!}{
\begin{tabular}{p{4cm} p{2cm} p{2cm}}
\toprule
\textbf{Configuration} & \textbf{Macro F1} & \textbf{Accuracy} \\
\midrule
Run A: Base Multimodal & 0.692 & 0.761 \\
Run B: Base + FGM & 0.724 & 0.793 \\
Run C: Full (FGM + DIPS + Stack) & \textbf{0.745} & \textbf{0.812} \\
\bottomrule
\end{tabular}
}
\end{table}

\subsection{Phase 2: CrisisMMD v2.0 Generalization}
To prove the architecture's generalizability across diverse linguistic and geographic boundaries, we applied the exact Run C framework to the CrisisMMD v2.0 `Damage Severity' task (using XLM-RoBERTa-Large). 

% -----------------------------
% Table 2
% -----------------------------
\begin{table}[h]
\centering
\caption{Generalization Benchmark on CrisisMMD (Phase 2)}
\label{table2}
\renewcommand{\arraystretch}{1.2}
\setlength{\tabcolsep}{6pt}
\resizebox{\columnwidth}{!}{
\begin{tabular}{p{4cm} p{2cm} p{2cm}}
\toprule
\textbf{Dataset Split} & \textbf{Macro F1} & \textbf{Accuracy} \\
\midrule
Hurricane Harvey & 0.751 & 0.820 \\
Hurricane Irma & 0.748 & 0.817 \\
\midrule
\textbf{Global Average} & \textbf{0.749} & \textbf{0.818} \\
\bottomrule
\end{tabular}
}
\end{table}

The results in Table \ref{table2} confirm that our noise-resilient fusion architecture, initially designed for low-resource Bengali pipelines, adapts seamlessly to high-resource multilingual disaster scenarios, maintaining a consistent Macro F1 around 0.75.

\section{Conclusion}
\fontsize{11}{13}\selectfont

This paper presents a robust, resource-constrained multimodal architecture for disaster severity assessment. By combining late-fusion with targeted FGM, safe FP16 training, DIPS, and Optuna-calibrated stacking, we effectively bridge the performance gap between low-resource (Bengali) and high-resource (Global) environments. Our open-sourced scripts provide a reproducible framework for deploying these life-saving models on edge and low-compute hardware during critical disaster responses.

\vspace{0.5cm}

\small
\textbf{Conflict of Interest:} The authors declare no conflict of interest.

\textbf{Author Contributions:} First Author contributed to conceptualization, methodology, and writing. Second Author and Third Author contributed to data analysis and review. All authors approved the final manuscript.

\textbf{Funding:} This research received no external funding.

\textbf{Ethical Statement:} This study follows ethical and academic integrity guidelines.

\balance
\fontsize{11}{13}\selectfont
\printbibliography

\end{document}
